\documentclass[11pt,a4paper]{article}

% ── Packages ──
\usepackage[utf8]{inputenc}
\usepackage[T1]{fontenc}
\usepackage{amsmath,amssymb}
\usepackage{physics}
\usepackage{siunitx}
\usepackage{booktabs}
\usepackage{geometry}
\usepackage{hyperref}
\usepackage{enumitem}
\usepackage{float}

\geometry{margin=2.5cm}
\hypersetup{colorlinks=true,linkcolor=blue,citecolor=blue,urlcolor=blue}

\title{Residual Gas Velocity Above a Stopped Rotating Disk\\
       in Dense Xenon: A Mechanism-by-Mechanism Analysis}
\author{}
\date{}

\begin{document}
\maketitle

% ══════════════════════════════════════════════════════════════
\section{Problem Definition}
\label{sec:problem}
% ══════════════════════════════════════════════════════════════

A solid horizontal disk of radius $R$ rotates inside a vertical
cylindrical vessel filled with xenon gas at high pressure. The disk
sweeps through a prescribed angle $\theta_\mathrm{sweep}$ using
bang-bang kinematics (maximum acceleration then maximum deceleration),
starting from rest and returning to rest. We seek the residual gas
velocity at a height $z_\mathrm{obs}$ above the disk surface, evaluated
at a time $t_\mathrm{obs}$ after the disk has stopped.

Unlike the propeller geometry (isolated NACA\,0012 blades), the disk
presents a \emph{continuous solid surface} of area $\pi R^2$ to the gas
above. The boundary layer is a von K\'arm\'an rotating-disk boundary
layer, which is qualitatively different from a flat-plate airfoil BL.

\subsection{Parameters}

Seven parameters are treated as variable inputs; all others are fixed.

\begin{table}[H]
\centering
\caption{Variable parameters (with default values).}
\label{tab:params}
\begin{tabular}{llll}
\toprule
Symbol & Description & Default & Units \\
\midrule
$h_\mathrm{disk}$        & Disk thickness                  & 0.005 & m   \\
$\theta_\mathrm{sweep}$  & Sweep angle                     & $\pi/8$ & rad \\
$m_\mathrm{disk}$        & Mass of disk                    & 50    & kg  \\
$m_\mathrm{plate}$       & Mass of each carrier plate      & 0.15  & kg  \\
$\tau$                    & Motor torque                    & 60    & N$\cdot$m \\
$z_\mathrm{obs}$         & Observation height above disk   & 0.005 & m \\
$t_\mathrm{obs}$         & Observation time after stop     & 0.5   & s   \\
\bottomrule
\end{tabular}
\end{table}

\begin{table}[H]
\centering
\caption{Fixed parameters.}
\label{tab:fixed}
\begin{tabular}{lll}
\toprule
Parameter & Value & Description \\
\midrule
$R$            & \SI{1.6}{m}        & Disk radius \\
$P$            & \SI{15}{bar}       & Xenon pressure \\
$T$            & \SI{300}{K}        & Gas temperature \\
$\Delta D$     & \SI{0.010}{m}      & Vessel inner diameter surplus \\
$H$            & \SI{1.5}{m}        & Vessel height \\
$N$            & 8                  & Number of carrier positions \\
\bottomrule
\end{tabular}
\end{table}

\subsection{Geometry and Coordinate System}

The disk rotates in the $xy$-plane with its upper surface at
$z = h_\mathrm{disk}$ above the vessel floor. The $z$-axis points
upward. The cylindrical vessel has inner radius
\begin{equation}
  R_\mathrm{cyl} = R + \frac{\Delta D}{2}\,,
\end{equation}
giving a radial gap $\Delta r = \Delta D/2 = \SI{5}{mm}$.
The gas column above the disk has height
$H_\mathrm{gas} = H - h_\mathrm{disk} \approx \SI{1.495}{m}$.


% ══════════════════════════════════════════════════════════════
\section{Xenon Gas Properties}
\label{sec:gas}
% ══════════════════════════════════════════════════════════════

Identical to the propeller analysis. At $P = \SI{15}{bar}$,
$T = \SI{300}{K}$:
\begin{alignat}{2}
  \rho  &= \frac{PM}{R_g T}
        &&\approx \SI{79}{kg/m^3}\,,
  \label{eq:rho}
  \\[4pt]
  \mu   &
        &&\approx \SI{2.32e-5}{Pa\cdot s}\,,
  \\[4pt]
  \nu   &= \mu/\rho
        &&\approx \SI{2.94e-7}{m^2/s}\,,
  \\[4pt]
  c_s   &= \sqrt{\gamma R_g T / M}
        &&\approx \SI{178}{m/s}\,.
\end{alignat}
The key feature: $\nu$ is \emph{very} small ($51\times$ smaller than
air at STP), leading to high Reynolds numbers at modest velocities and
extremely slow molecular diffusion.


% ══════════════════════════════════════════════════════════════
\section{Kinematics}
\label{sec:kinematics}
% ══════════════════════════════════════════════════════════════

\subsection{Moment of Inertia}

The disk is a uniform solid cylinder of mass $m_\mathrm{disk}$ and
radius $R$:
\begin{equation}
  I_\mathrm{disk} = \tfrac{1}{2}\,m_\mathrm{disk}\,R^2\,.
  \label{eq:Idisk}
\end{equation}
The $N$ carrier plates are housed in the hangar at $r \approx R$:
\begin{equation}
  I_\mathrm{plates} = N\,m_\mathrm{plate}\,R^2\,.
  \label{eq:Iplates}
\end{equation}
Total:
\begin{equation}
  I = I_\mathrm{disk} + I_\mathrm{plates}
    = \tfrac{1}{2}\,m_\mathrm{disk}\,R^2 + N\,m_\mathrm{plate}\,R^2\,.
  \label{eq:Itotal}
\end{equation}

For default parameters ($m_\mathrm{disk} = \SI{50}{kg}$,
$m_\mathrm{plate} = \SI{0.15}{kg}$, $N = 8$):
$I_\mathrm{disk} = \SI{64.0}{kg\cdot m^2}$,
$I_\mathrm{plates} = \SI{3.07}{kg\cdot m^2}$,
$I = \SI{67.1}{kg\cdot m^2}$.

\subsection{Aerodynamic Drag Torque}

A rotating disk in a fluid experiences a resisting torque from skin
friction on both the upper and lower surfaces. For a turbulent boundary
layer on one side of the disk, the moment coefficient is
(Schlichting \& Gersten, \emph{Boundary-Layer Theory}, 9th ed., Ch.~5):
\begin{equation}
  C_M = \frac{0.146}{\mathrm{Re}_R^{1/5}}\,,
  \label{eq:CM}
\end{equation}
where $\mathrm{Re}_R = \omega R^2 / \nu$ is the rotational Reynolds
number. The torque on one side is
\begin{equation}
  \tau_\mathrm{one\,side} = \tfrac{1}{2}\,C_M\,\rho\,\omega^2\,R^5\,.
  \label{eq:tau_side}
\end{equation}
For two sides (upper and lower):
\begin{equation}
  \tau_\mathrm{drag} = C_M\,\rho\,\omega^2\,R^5\,.
  \label{eq:tau_drag}
\end{equation}
Note that $C_M$ depends on $\omega$ through $\mathrm{Re}_R$, making
the drag nonlinear. For the edge contribution (rim drag), we add a
bluff-body drag on the annular edge of thickness $h_\mathrm{disk}$:
\begin{equation}
  \tau_\mathrm{rim} = \tfrac{1}{2}\,\rho\,C_{d,\mathrm{rim}}\,
    h_\mathrm{disk}\,R^4\,\omega^2\,,
  \label{eq:tau_rim}
\end{equation}
with $C_{d,\mathrm{rim}} \approx 1.0$ (blunt edge, Hoerner,
\emph{Fluid-Dynamic Drag}, Ch.~3).

\subsection{Bang-Bang Motion}

The equation of motion is
\begin{equation}
  I\,\frac{d\omega}{dt} = \pm\tau_\mathrm{motor}
    - \tau_\mathrm{drag}(\omega) - \tau_\mathrm{rim}(\omega)\,,
\end{equation}
where the sign alternates between acceleration and deceleration phases.
The switching angle $\theta_\mathrm{switch}$ is found numerically such
that $\omega \to 0$ exactly when $\theta \to \theta_\mathrm{sweep}$.


% ══════════════════════════════════════════════════════════════
\section{Boundary Layer: Von K\'arm\'an Rotating Disk}
\label{sec:BL}
% ══════════════════════════════════════════════════════════════

\subsection{The Classical Problem}

The boundary layer on an infinite rotating disk in an otherwise
quiescent fluid was solved by von K\'arm\'an (1921) and refined by
Cochran (1934). The key result: the BL thickness is
\emph{independent of radius} and depends only on rotation rate and
viscosity:
\begin{equation}
  \delta = a\,\sqrt{\frac{\nu}{\omega}}\,,
  \label{eq:delta_vK}
\end{equation}
where $a \approx 5.5$ (laminar, Schlichting \& Gersten, Ch.~5.2).
The BL has three velocity components:
\begin{itemize}
  \item \textbf{Tangential} ($v_\phi$): the gas near the disk is
        dragged in the direction of rotation.
  \item \textbf{Radial} ($v_r$): centrifugal force pushes the gas
        outward within the BL.
  \item \textbf{Axial} ($v_z$): to replace the gas expelled radially,
        fresh gas is drawn downward from above the BL. This is the
        \emph{von K\'arm\'an pumping} effect.
\end{itemize}

The axial (pumping) velocity far from the disk is
(von K\'arm\'an, 1921; Cochran, 1934):
\begin{equation}
  w_\infty = -0.886\,\sqrt{\nu\,\omega}
  \qquad\text{(laminar)}\,.
  \label{eq:w_inf_lam}
\end{equation}
The negative sign indicates flow \emph{toward} the disk. This means
the rotating disk \emph{sucks gas down from above} uniformly across
its entire surface.

\subsection{Turbulent Regime}

The rotating-disk BL transitions to turbulence at
$\mathrm{Re}_r = \omega r^2 / \nu \approx 2.5 \times 10^5$
(Lingwood, 1995). In the turbulent regime, the BL thickness grows:
\begin{equation}
  \delta_\mathrm{turb} \approx 0.526\,r\,
    \left(\frac{\nu}{\omega r^2}\right)^{1/5}
  = 0.526\,r\,\mathrm{Re}_r^{-1/5}\,,
  \label{eq:delta_turb}
\end{equation}
(Owen \& Rogers, \emph{Flow and Heat Transfer in Rotating-Disc Systems},
Vol.~1, 1989). Unlike the laminar case, the turbulent BL thickness
\emph{increases with radius}.

The turbulent axial pumping velocity also increases:
\begin{equation}
  w_\infty^\mathrm{turb} \approx -C_w\,\sqrt{\nu_t\,\omega}\,,
  \label{eq:w_inf_turb}
\end{equation}
where the turbulent viscosity $\nu_t \gg \nu$. A more practical
estimate uses the entrainment flow rate. The volumetric flow entrained
by one side of the disk (turbulent) is
(Owen \& Rogers, 1989):
\begin{equation}
  Q = 0.219\,\mathrm{Re}_R^{4/5}\,\nu\,R\,,
  \label{eq:Q_turb}
\end{equation}
where $\mathrm{Re}_R = \omega R^2 / \nu$. The mean axial velocity at
the disk edge (over the disk area):
\begin{equation}
  \bar{w}_\mathrm{axial} = \frac{Q}{\pi R^2}\,.
  \label{eq:w_mean}
\end{equation}

\subsection{Relevance to the MARS Problem}

This is the fundamental difference from the propeller: a rotating
disk has an \emph{intrinsic axial pumping velocity} that draws gas
\emph{toward} the disk surface. During rotation, this sets up a
meridional circulation (radially outward in the BL, axially downward
above the disk). After the disk stops, this circulation must decay.

The question is: how quickly does the axial flow decay, and what is its
magnitude at $z = z_\mathrm{obs}$ at time $t = t_\mathrm{obs}$?


% ══════════════════════════════════════════════════════════════
\section{Mechanism 1: Von K\'arm\'an Pumping and Decay}
\label{sec:pumping}
% ══════════════════════════════════════════════════════════════

\subsection{Axial Velocity During Rotation}

During steady rotation at angular velocity $\omega$, the axial
(pumping) velocity depends on height above the disk. In the laminar
regime (Cochran's solution):
\begin{equation}
  w(z) = -\sqrt{\nu\omega}\;H_\mathrm{vK}(\zeta)\,,
  \qquad \zeta = z\sqrt{\omega/\nu}\,,
\end{equation}
where $H_\mathrm{vK}(\zeta)$ is a universal function that increases
from $0$ at $\zeta = 0$ (disk surface, no-penetration) to $0.886$ as
$\zeta \to \infty$. The asymptotic value is reached for
$\zeta \gtrsim 5$ (i.e.\ $z \gtrsim \delta$).

For the turbulent regime at the outer part of the disk, the pumping is
stronger. The mean downward velocity can be estimated from the
entrainment flow rate~\eqref{eq:Q_turb}.

\subsection{Decay After Stopping}

When the disk stops at $t = 0$, the forcing that sustains the BL and
the pumping vanishes instantaneously. The decay proceeds through
several stages:

\paragraph{Stage 1: Turbulent decay ($0 < t \lesssim 2\,t_e$).}
The turbulent eddies in the BL lose their energy source and decay on
the eddy turnover timescale:
\begin{equation}
  t_e = \frac{\delta}{\omega\,r \cdot 0.05}\,,
\end{equation}
where $0.05\,\omega r$ is the turbulence intensity. The turbulent
viscosity decays exponentially with $e$-folding time $\sim 2\,t_e$.
During this period the BL profile changes from a turbulent
$1/7$-power law to a lower-energy state.

\paragraph{Stage 2: Molecular spin-down ($t \gg 2\,t_e$).}
After the turbulence dies, only molecular viscosity remains. The
tangential and radial velocity components diffuse outward from the
disk surface. The tangential velocity profile evolves as
\begin{equation}
  v_\phi(z,t) \sim V_\phi(0)\,
    \mathrm{erfc}\!\left(\frac{z}{2\sqrt{\nu t}}\right),
\end{equation}
where $V_\phi(0)$ is the initial tangential velocity at the edge of
the decayed BL. This is the classical Stokes first problem
(Schlichting \& Gersten, Ch.~4.3).

\paragraph{Axial velocity.}
The axial pumping is \emph{driven by} the tangential and radial
velocities in the BL via the continuity equation. Once the tangential
velocity decays (no more centrifugal force in the BL), the radial
outflow ceases, and the axial inflow ceases with it. The axial
velocity decays on the \emph{same timescale} as the tangential BL.

The axial pumping velocity scales as $w \sim \sqrt{\nu_\mathrm{eff}\,
\omega_\mathrm{eff}}$ where $\omega_\mathrm{eff}$ is the effective
angular velocity of the gas near the disk surface. As
$\omega_\mathrm{eff} \to 0$, $w \to 0$.

\subsection{Quantitative Estimate}

\paragraph{During rotation (peak values).}
The peak angular velocity $\omega_\mathrm{max}$ is obtained from the
bang-bang kinematics. The laminar pumping velocity:
\begin{equation}
  w_\mathrm{lam} = 0.886\,\sqrt{\nu\,\omega_\mathrm{max}}\,.
  \label{eq:w_lam_peak}
\end{equation}

The turbulent entrainment flow rate gives the mean axial velocity
over the disk area. This represents the steady-state value during
rotation.

\paragraph{After stop.}
The PDE governing the tangential velocity is
\begin{equation}
  \frac{\partial v_\phi}{\partial t}
  = \frac{\partial}{\partial z}\left[
    \nu_\mathrm{eff}(z,t)\,\frac{\partial v_\phi}{\partial z}
  \right],
  \label{eq:vphi_PDE}
\end{equation}
with $v_\phi(0,t) = 0$ (no-slip on stopped disk),
$v_\phi(z \to \infty, t) = 0$, and initial condition from the
von K\'arm\'an profile. This is essentially the same 1D diffusion
equation as for the propeller, but with a different initial condition
and with the crucial difference that the initial profile extends over
the entire disk area, not just $N$ narrow blade wakes.

The axial velocity at $z_\mathrm{obs}$ is then obtained from continuity:
\begin{equation}
  w(z_\mathrm{obs}, t) \approx -\frac{\partial}{\partial r}
    \int_0^{z_\mathrm{obs}} v_r(r,z,t)\,dz\,.
\end{equation}
In practice, for $z_\mathrm{obs} > \delta$, the axial velocity is
estimated as $w \sim \sqrt{\nu_\mathrm{eff}(t) \cdot
\omega_\mathrm{eff}(t)}$, which decays as both $\nu_t$ and
$\omega_\mathrm{eff}$ decay.

\subsection{Result}

The von K\'arm\'an pumping produces a nonzero axial velocity
\emph{during} rotation. After the disk stops, this axial velocity
decays on the same timescale as the tangential BL. For
$t_\mathrm{obs} \gg 2\,t_e$, the pumping has ceased. However, the
magnitude and decay time depend strongly on the disk angular velocity.
The companion Python script evaluates this numerically.

\begin{equation}
  \boxed{u(z_\mathrm{obs},\, t_\mathrm{obs}) = \text{computed numerically;
  see Python script}}
\end{equation}


% ══════════════════════════════════════════════════════════════
\section{Mechanism 2: Tangential BL Diffusion}
\label{sec:tangential}
% ══════════════════════════════════════════════════════════════

Even without the axial pumping, the tangential (azimuthal) velocity
in the BL can in principle diffuse vertically, exactly as in the
propeller case.

\subsection{Initial Tangential Profile}

The von K\'arm\'an tangential velocity profile (turbulent,
$1/7$-power law approximation):
\begin{equation}
  v_\phi(z) = \omega r \left[1 - \left(\frac{z}{\delta}\right)^{1/7}
    \right]
  \qquad (0 < z \le \delta)\,,
\end{equation}
where $\delta$ is the local turbulent BL thickness. At the tip
($r = R$), $v_\phi(0) = \omega R = V_\mathrm{tip}$.

\subsection{Diffusion Equation}

Same 1D diffusion equation as for the propeller:
\begin{equation}
  \frac{\partial v_\phi}{\partial t}
  = \frac{\partial}{\partial z}
    \left[\nu_\mathrm{eff}(z,t)\,\frac{\partial v_\phi}{\partial z}
    \right],
\end{equation}
with initial condition from the von K\'arm\'an profile and decaying
$\nu_t$.

\subsection{Key Difference from Propeller}

The propeller BL has thickness $\delta = 0.37\,c\,\mathrm{Re}_c^{-1/5}$,
which depends on the chord $c$. The disk BL thickness at radius $r$ is
$\delta = 0.526\,r\,\mathrm{Re}_r^{-1/5}$, which depends on \emph{radius}.
At the outer edge ($r = R = \SI{1.6}{m}$), the disk BL is much thicker
than a propeller blade BL (because $R \gg c$).

Whether $z_\mathrm{obs} - \delta > 0$ (observation above the BL) or
$z_\mathrm{obs} - \delta < 0$ (observation inside the BL) depends on
the disk angular velocity and observation height.

\subsection{Result}

The 1D PDE solver evaluates the tangential velocity at $z_\mathrm{obs}$,
$t_\mathrm{obs}$. If $z_\mathrm{obs} > \delta + \ell_\mathrm{turb}
+ \ell_\mathrm{mol}$, the result is zero. If the observation point is
close to $\delta$, a small nonzero tail may exist (same as propeller).

\begin{equation}
  \boxed{v_\phi(z_\mathrm{obs},\, t_\mathrm{obs}) = \text{computed
  numerically; see Python script}}
\end{equation}


% ══════════════════════════════════════════════════════════════
\section{Mechanism 3: Bulk Swirl / Ekman Spin-Down}
\label{sec:swirl}
% ══════════════════════════════════════════════════════════════

\subsection{Angular Impulse}

The disk drag torque deposits angular momentum in the gas. Unlike the
propeller (where the drag acts on $N$ narrow blades), the disk torque
acts over the entire disk surface:
\begin{equation}
  \tau_\mathrm{drag}
    = C_M\,\rho\,\omega^2\,R^5 + \tfrac{1}{2}\,\rho\,C_{d,\mathrm{rim}}
      \,h_\mathrm{disk}\,R^4\,\omega^2\,.
\end{equation}
The angular impulse over the sweep is $J = \int_0^{T_\mathrm{sweep}}
\tau_\mathrm{drag}\,dt$.

\subsection{Ekman Spin-Down}

In a confined vessel, the bulk swirl decays via the Ekman mechanism:
endcap boundary layers transport angular momentum to the walls. The
Ekman spin-down time is
\begin{equation}
  t_\mathrm{Ekman} = \frac{H}{\sqrt{\nu\,\Omega_\mathrm{gas}}}\,,
\end{equation}
where $\Omega_\mathrm{gas}$ is the characteristic gas rotation rate.
For molecular $\nu$, this is $\sim\SI{1000}{s}$; for turbulent $\nu_t$
(lasting $\sim\SI{10}{ms}$), $\sim\SI{50}{s}$.

The Ekman pumping velocity $w_E = \sqrt{\nu\,\Omega}$ is of order
$\sim\SI{1}{mm/s}$ (molecular), producing negligible vertical
displacement in $t_\mathrm{obs}$.

\subsection{Result}

\begin{equation}
  \boxed{u(z_\mathrm{obs},\, t_\mathrm{obs})
    \approx 0 \quad \text{(bulk Ekman spin-down)}}
\end{equation}
Same conclusion as for the propeller: the Ekman mechanism is far too
slow to affect the observation.


% ══════════════════════════════════════════════════════════════
\section{Mechanism 4: Displacement (Potential) Flow}
\label{sec:potential}
% ══════════════════════════════════════════════════════════════

As the disk rotates, the rim (a blunt edge of height $h_\mathrm{disk}$)
pushes gas aside. The displacement flow around the rim produces a
transient velocity field that decays at the speed of sound after the
disk stops.

\subsection{Clearing Time}

\begin{equation}
  t_\mathrm{acoustic} = \frac{R}{c_s}
    \approx \frac{\SI{1.6}{m}}{\SI{178}{m/s}}
    \approx \SI{9}{ms}\,.
\end{equation}

\subsection{Blockage}

The disk blockage ratio is $h_\mathrm{disk}/\Delta r = 1.0$ (for
default parameters), much smaller than the propeller blade blockage
(5.3). The displacement flow is correspondingly weaker.

\subsection{Result}

\begin{equation}
  \boxed{u(z_\mathrm{obs},\, t_\mathrm{obs})
    = 0 \quad \text{(displacement flow, for }
    t_\mathrm{obs} \gg t_\mathrm{acoustic}\text{)}}
\end{equation}


% ══════════════════════════════════════════════════════════════
\section{Mechanism 5: Pressure Pulse from Deceleration}
\label{sec:pressure}
% ══════════════════════════════════════════════════════════════

Same physics as for the propeller. The deceleration of the disk
creates a transient pressure perturbation:
\begin{equation}
  \Delta p \sim \rho\,V_\mathrm{tip}^2\,,
\end{equation}
which propagates as an acoustic wave passing through $z_\mathrm{obs}$
in $\sim z_\mathrm{obs}/c_s$ and leaving no residual.

\begin{equation}
  \boxed{u(z_\mathrm{obs},\, t_\mathrm{obs})
    = 0 \quad \text{(pressure pulse, for }
    t_\mathrm{obs} \gg t_\mathrm{decel}\text{)}}
\end{equation}


% ══════════════════════════════════════════════════════════════
\section{Mechanism 6: Rim Separation and Wake}
\label{sec:rim}
% ══════════════════════════════════════════════════════════════

The disk has a blunt rim of height $h_\mathrm{disk}$, which produces
flow separation at the edge during rotation. Unlike the streamlined
NACA\,0012 trailing edge of the propeller, the disk edge generates a
separated wake.

\subsection{Wake Character}

The separated wake at the rim is \emph{tangential}: it lies in the
plane of the disk edge, in the annular gap between the disk and the
vessel wall. The wake has a vertical extent of order $\sim h_\mathrm{disk}
= \SI{5}{mm}$, and its velocity decays after the disk stops.

\subsection{Vertical Transport}

The rim wake is tangential and confined to the disk-edge plane. It does
not have an axial velocity component directed upward toward
$z_\mathrm{obs}$. The same transport limitations as for the propeller
vortex shedding apply: vertical diffusion from the rim wake is limited
by the same BL-thickness-plus-diffusion reach constraint.

\subsection{Result}

\begin{equation}
  \boxed{u(z_\mathrm{obs},\, t_\mathrm{obs})
    = 0 \quad \text{(rim wake, confined to edge)}}
\end{equation}


% ══════════════════════════════════════════════════════════════
\section{Mechanism 7: Secondary Flows (Centrifugal, Coriolis)}
\label{sec:secondary}
% ══════════════════════════════════════════════════════════════

\subsection{Centrifugal Effects in the Bulk}

Gas in the bulk (above the BL) that has acquired angular momentum is
subject to centrifugal force. This drives radial outflow, which must
be compensated by axial inflow from above (this is the von K\'arm\'an
pumping, already treated in Section~\ref{sec:pumping}).

After the disk stops, the centrifugal force decays with the angular
velocity of the gas. No new mechanism beyond what is already accounted
for in Section~\ref{sec:pumping}.

\subsection{Rim-Gap Couette Flow}

In the radial clearance $\Delta r$ between the disk edge and the vessel
wall:
\begin{equation}
  \mathrm{Re}_\mathrm{gap}
    = \frac{V_\mathrm{tip}\,\Delta r}{\nu}\,.
\end{equation}
This flow is purely tangential (azimuthal shear between moving disk
edge and stationary wall), with no axial velocity component. Same
conclusion as for the propeller tip gap.

\subsection{Result}

\begin{equation}
  \boxed{u(z_\mathrm{obs},\, t_\mathrm{obs})
    = 0 \quad \text{(secondary flows)}}
\end{equation}


% ══════════════════════════════════════════════════════════════
\section{Vessel Confinement Effects}
\label{sec:confinement}
% ══════════════════════════════════════════════════════════════

The same confinement analysis as for the propeller applies. All five
effects (wall BLs, Ekman spin-up, tip-gap Couette, reflected pressure
waves, displacement readjustment) are quantified in the companion
Python script and shown to be negligible. See the propeller analysis
(\texttt{propeller\_flow\_analysis.tex}, Section~7) for the detailed
treatment; the physics is identical.


% ══════════════════════════════════════════════════════════════
\section{Key Difference from Propeller: Axial Pumping}
\label{sec:comparison}
% ══════════════════════════════════════════════════════════════

The propeller analysis found that \emph{all} mechanisms give zero at
$z_\mathrm{obs}$ (provided $z_\mathrm{obs} - \delta > \sim\SI{2}{mm}$).
For the rotating disk, the situation is qualitatively different because
of the \textbf{von K\'arm\'an pumping}:

\begin{itemize}
  \item The propeller blades have no axial flow: at $\alpha = 0$, the
        NACA\,0012 produces no lift and no axial pressure gradient. All
        flow is tangential.
  \item The rotating disk intrinsically generates an axial velocity
        ($w_\infty = -0.886\sqrt{\nu\omega}$ in the laminar regime,
        stronger in the turbulent regime). This is a \emph{volume effect}:
        the entire disk area participates.
  \item The pumping extends to $z \to \infty$ (in an infinite domain),
        meaning it reaches $z_\mathrm{obs}$ during rotation.
  \item After the disk stops, the pumping decays as the BL tangential
        velocity decays. The question is whether this decay is fast
        enough.
\end{itemize}

The BL thickness is also much larger for the disk than for the
propeller. At the disk edge ($r = R = \SI{1.6}{m}$):
\begin{equation}
  \delta_\mathrm{disk}(R) = 0.526\,R\,\mathrm{Re}_R^{-1/5}\,,
\end{equation}
compared with the propeller:
\begin{equation}
  \delta_\mathrm{prop} = 0.37\,c\,\mathrm{Re}_c^{-1/5}\,.
\end{equation}
Since $R = \SI{1.6}{m} \gg c = \SI{0.22}{m}$, the disk BL is
significantly thicker, potentially engulfing $z_\mathrm{obs}$.

Both quantities (pumping velocity and BL thickness) are evaluated
numerically in the companion Python script.


% ══════════════════════════════════════════════════════════════
\section{Conclusion}
\label{sec:conclusion}
% ══════════════════════════════════════════════════════════════

Table~\ref{tab:summary} summarises all mechanisms.

\begin{table}[H]
\centering
\caption{Summary of all mechanisms for residual velocity at
$z_\mathrm{obs}$ above the disk, $t_\mathrm{obs}$ after stopping.}
\label{tab:summary}
\begin{tabular}{llc}
\toprule
Mechanism & Characteristic timescale / Comment
  & $u(z_\mathrm{obs}, t_\mathrm{obs})$ \\
\midrule
Von K\'arm\'an pumping (axial)
  & Decays with BL tangential velocity
  & Numerical \\
Tangential BL diffusion
  & Same as propeller eddy diffusion
  & Numerical \\
Bulk Ekman spin-down
  & $t_\mathrm{Ekman} \sim \SI{1000}{s}$ (mol.)
  & 0 \\
Displacement flow
  & $t_\mathrm{acoustic} \approx \SI{9}{ms}$
  & 0 \\
Pressure pulse
  & $t_\mathrm{decel} \sim \SI{10}{ms}$
  & 0 \\
Rim wake
  & Confined to disk edge
  & 0 \\
Secondary flows
  & Centrifugal decays with $\omega$; gap flow tangential
  & 0 \\
\bottomrule
\end{tabular}
\end{table}

The critical difference from the propeller is the \textbf{von K\'arm\'an
pumping}, which provides a direct axial velocity mechanism that extends
above the disk BL. The propeller has no such mechanism. Whether this
makes the disk geometry worse depends on the decay rate of the pumping
after the disk stops, which is evaluated quantitatively in the companion
script \texttt{disk\_flow\_numerics.py}.

A second concern is the \textbf{BL thickness}: the disk BL at $r = R$
may be thick enough to engulf $z_\mathrm{obs}$, in which case the
observation point is inside the BL and the velocity is necessarily
nonzero. The propeller BL (chord-based) is thinner.

\end{document}
